Calcul de l'auteur. Variation de la taxe annuelle de chaque bien : bascule à recette constante vers une taxe sur le terrain (figure précédente). Prix : variation de la taxe rapportée à la valeur du bien, multipliée par la semi-élasticité estimée (4,26 \% de prix par point pour mille chez Høj et al. ; 2,7 \% pour une capitalisation complète au taux de 3,7 \% de Koster et Pinchbeck ; 1,25 \% au plus chez Nielsson et al., dont l'estimation centrale est nulle). Loyers : part de la variation de taxe répercutée sur le loyer.\par Lecture : si l'élasticité de Høj et al. valait en France, l'appartement parisien perdrait 12 \% de sa valeur ; avec celle de Nielsson et al., au plus 4 \%. Les effets danois et anglais portent sur des variations de taxe de quelques points pour mille ; la baisse de taxe de la maison corrézienne (8 ‰ de sa valeur) sort de ce domaine.
